\documentclass[10pt,a4paper]{article}

\usepackage[T1]{fontenc}
\usepackage[utf8]{inputenc}
\usepackage{lmodern}
\usepackage{microtype}
\usepackage[a4paper,top=23mm,bottom=21mm,left=23mm,right=23mm]{geometry}
\usepackage{amsmath,amssymb,mathtools}
\usepackage{booktabs,array,tabularx}
\usepackage{enumitem}
\usepackage[table]{xcolor}
\usepackage{tikz}
\usepackage[most]{tcolorbox}
\usepackage{fancyhdr}
\usepackage{titlesec}
\usepackage{hyperref}

\definecolor{Navy}{HTML}{14253D}
\definecolor{Teal}{HTML}{178C8C}
\definecolor{PaleTeal}{HTML}{EAF6F5}
\definecolor{Gold}{HTML}{F2B84B}
\definecolor{Ink}{HTML}{243447}
\definecolor{Muted}{HTML}{66788A}
\definecolor{Paper}{HTML}{F7F9FB}
\definecolor{Line}{HTML}{D8E1E8}

\hypersetup{
  colorlinks=true,
  linkcolor=Teal,
  urlcolor=Teal,
  pdftitle={Diagnostica tensoriale irriducibile per un sistema spin-j},
  pdfsubject={Base Tlm e validazione del decadimento multipolare},
  pdfauthor={Codex}
}

\color{Ink}
\setlength{\parindent}{0pt}
\setlength{\parskip}{4.5pt}
\setlist[itemize]{leftmargin=5mm,itemsep=2pt,topsep=3pt}
\renewcommand{\arraystretch}{1.25}

\pagestyle{fancy}
\setlength{\headheight}{21pt}
\fancyhf{}
\fancyhead[L]{\sffamily\scriptsize\color{Muted}\MakeUppercase{Spin quantum diffusion - nota tecnica}}
\fancyhead[R]{\color{Line}\rule{25mm}{0.35pt}}
\fancyfoot[L]{\tikz\fill[Teal] (0,0) circle (1.25mm);}
\fancyfoot[R]{\sffamily\scriptsize\color{Muted}\thepage}
\renewcommand{\headrulewidth}{0.35pt}
\renewcommand{\headrule}{\hbox to\headwidth{\color{Line}\leaders\hrule height \headrulewidth\hfill}}

\titleformat{\section}
  {\sffamily\bfseries\fontsize{19}{22}\selectfont\color{Navy}}
  {}{0pt}{}
\titlespacing*{\section}{0pt}{0pt}{8pt}

\titleformat{\subsection}
  {\sffamily\bfseries\large\color{Teal}}
  {}{0pt}{}
\titlespacing*{\subsection}{0pt}{10pt}{3pt}

\newcolumntype{Y}{>{\raggedright\arraybackslash}X}
\newcolumntype{C}[1]{>{\centering\arraybackslash}p{#1}}

\newtcolorbox{equationbox}{
  enhanced,
  colback=Paper,
  colframe=Line,
  boxrule=0.45pt,
  arc=0pt,
  left=2mm,right=2mm,top=1mm,bottom=1mm,
  before skip=5pt,after skip=6pt
}

\newtcolorbox{notebox}[1][]{
  enhanced,
  colback=white,
  colframe=Line,
  boxrule=0.5pt,
  borderline west={2.4pt}{0pt}{Teal},
  arc=0pt,
  left=3.5mm,right=3mm,top=2.5mm,bottom=2.5mm,
  before upper=\raggedright,
  #1
}

\newcommand{\cardtitle}[1]{{\sffamily\bfseries\color{Navy}#1\par}\smallskip}
\newcommand{\muted}[1]{{\color{Muted}#1}}
\newcommand{\HS}{\mathrm{HS}}

\begin{document}

% -----------------------------------------------------------------------------
% Copertina
% -----------------------------------------------------------------------------
\thispagestyle{fancy}
\vspace*{22mm}

{\sffamily\bfseries\small\color{Teal}\MakeUppercase{Nota esplicativa}}\par
\vspace{5mm}
{\sffamily\bfseries\fontsize{27}{31}\selectfont\color{Navy}
Diagnostica tensoriale\\[1mm]
irriducibile per uno spin-$j$\par}

\vspace{5mm}
{\color{Gold}\rule{22mm}{2.2mm}}

\vspace{9mm}
{\fontsize{12}{17}\selectfont\color{Muted}
Come leggere i tensori sferici $T_{\ell m}$, perch\'e vengono costruiti
algebricamente e in che modo verificano la legge di diffusione
$e^{-D\ell(\ell+1)t}$.\par}

\vspace{10mm}
\begin{tcbraster}[
  raster columns=3,
  raster equal height=rows,
  raster column skip=3mm
]
  \begin{notebox}
    \cardtitle{1 - Scomporre}
    La matrice densit\`a viene riscritta come somma di multipoli: monopolo,
    dipolo, quadrupolo e strutture di rango superiore.
  \end{notebox}
  \begin{notebox}[borderline west={2.4pt}{0pt}{Gold}]
    \cardtitle{2 - Evolvere}
    La diffusione isotropa agisce separatamente su ogni rango $\ell$ e non
    mescola multipoli differenti.
  \end{notebox}
  \begin{notebox}[borderline west={2.4pt}{0pt}{Navy}]
    \cardtitle{3 - Verificare}
    I coefficienti numerici vengono confrontati con il decadimento teorico
    previsto per ogni coppia $(\ell,m)$.
  \end{notebox}
\end{tcbraster}

\vspace{9mm}
\subsection{Idea centrale}
Il modulo non introduce una nuova dinamica fisica. Fornisce un sistema di
coordinate naturale per osservare ci\`o che la dinamica esistente sta facendo:
\`e una diagnostica spettrale dello stato di spin.

\begin{equationbox}
\[
  \rho(t)
  \;\longrightarrow\;
  \{c_{\ell m}(t)\}
  \;\longrightarrow\;
  c_{\ell m}(t)\stackrel{?}{=}e^{-D\ell(\ell+1)t}c_{\ell m}(0)
\]
\end{equationbox}

\vspace{5mm}
\muted{Contesto: sistema singolo di spin-$j$, dimensione $d=2j+1$. Il caso
predefinito del prototipo \`e $j=2$, quindi $d=5$.}

\newpage

% -----------------------------------------------------------------------------
% 1. Base multipolare
% -----------------------------------------------------------------------------
\section*{1. Una base multipolare per la matrice densit\`a}

Una matrice densit\`a di dimensione $d=2j+1$ appartiene a uno spazio di
operatori di dimensione $d^2$. I tensori sferici irriducibili
$T_{\ell m}$ formano una base completa di questo spazio:

\begin{equationbox}
\[
  \rho
  =\sum_{\ell=0}^{2j}\sum_{m=-\ell}^{\ell}
    c_{\ell m}T_{\ell m}.
\]
\end{equationbox}

Questo \`e l'analogo quantistico dell'espansione di una funzione sulla sfera
in armoniche sferiche $Y_{\ell m}$.

\subsection{Che cosa rappresenta il rango \texorpdfstring{$\ell$}{l}}

\rowcolors{2}{Paper}{white}
\begin{tabularx}{\textwidth}{>{\bfseries}p{20mm}YY}
  \rowcolor{Navy}
  \color{white}Rango & \color{white}\bfseries Interpretazione &
  \color{white}\bfseries Comportamento \\
  $\ell=0$ & Monopolo: traccia e normalizzazione & Invariato dalla diffusione \\
  $\ell=1$ & Dipolo: orientamento medio & Struttura angolare grossolana \\
  $\ell=2$ & Quadrupolo: anisotropia di secondo ordine & Decade pi\`u rapidamente del dipolo \\
  $\ell>2$ & Multipoli superiori & Dettagli angolari progressivamente pi\`u fini \\
\end{tabularx}

\subsection{Perch\'e si chiamano irriducibili}

Sotto una rotazione, gli operatori con lo stesso rango $\ell$ possono
mescolarsi attraverso l'indice $m$, ma non si mescolano con ranghi diversi.
Le relazioni caratteristiche sono

\begin{equationbox}
\[
 [J_z,T_{\ell m}]=mT_{\ell m},
 \qquad
 [J_{\pm},T_{\ell m}]
 =\sqrt{(\ell\mp m)(\ell\pm m+1)}\,T_{\ell,m\pm1}.
\]
\end{equationbox}

In altre parole, per ogni valore di $\ell$ i $2\ell+1$ componenti formano un
sottospazio chiuso rispetto alle rotazioni.

\begin{notebox}[colback=PaleTeal,colframe=Teal]
  \begin{tabularx}{\linewidth}{>{\bfseries}p{35mm}YY}
    Esempio: $j=2$ & $\ell=0,1,2,3,4$ &
    $1+3+5+7+9=25=d^2$
  \end{tabularx}
\end{notebox}

La somma delle dimensioni $2\ell+1$ restituisce l'intero spazio degli
operatori $5\times5$.

\newpage

% -----------------------------------------------------------------------------
% 2. Costruzione algebrica
% -----------------------------------------------------------------------------
\section*{2. Normalizzazione, fase e costruzione algebrica}

\subsection{Ortonormalit\`a di Hilbert-Schmidt}

Gli operatori vengono trattati come vettori mediante il prodotto scalare

\begin{equationbox}
\[
  \langle A,B\rangle_{\HS}=\operatorname{Tr}(A^{\dagger}B).
\]
\end{equationbox}

La base viene normalizzata in modo che

\begin{equationbox}
\[
  \operatorname{Tr}\!\left(
    T_{\ell m}^{\dagger}T_{\ell' m'}
  \right)
  =\delta_{\ell\ell'}\delta_{mm'}.
\]
\end{equationbox}

Di conseguenza, le coordinate della matrice densit\`a si leggono direttamente:

\begin{equationbox}
\[
  c_{\ell m}=\operatorname{Tr}(\rho T_{\ell m}^{\dagger}),
  \qquad
  P_{\ell}=\sum_{m=-\ell}^{\ell}|c_{\ell m}|^2.
\]
\end{equationbox}

$P_{\ell}$ misura quanta struttura dello stato risiede complessivamente nel
rango $\ell$.

\subsection{Convenzione di fase Condon-Shortley}

I tensori possono essere moltiplicati per fasi senza perdere ortogonalit\`a.
Per rendere i segni confrontabili con la convenzione standard viene scelta una
fase coerente. In questa implementazione compare il fattore $(-1)^{\ell}$ e si
ottiene

\begin{equationbox}
\[
  T_{\ell m}^{\dagger}=(-1)^mT_{\ell,-m}.
\]
\end{equationbox}

La fase non cambia $P_{\ell}$, ma conta quando si confrontano coefficienti
complessi individuali o implementazioni differenti.

\subsection{Dal peso massimo agli altri componenti}

Per ogni rango $\ell$ si parte dall'operatore di peso massimo
$J_+^{\ell}$, con $J_+=J_x+iJ_y$. Dopo la normalizzazione, gli altri valori di
$m$ vengono generati tramite il commutatore con $J_-=J_x-iJ_y$:

\begin{equationbox}
\begin{align*}
  T_{\ell\ell}
  &=(-1)^{\ell}
    \frac{J_+^{\ell}}{\lVert J_+^{\ell}\rVert_{\HS}},\\[2mm]
  T_{\ell,m-1}
  &=\frac{[J_-,T_{\ell m}]}
  {\sqrt{(\ell+m)(\ell-m+1)}}.
\end{align*}
\end{equationbox}

Qui il commutatore non abbassa uno stato $\lvert j,m\rangle$: abbassa l'indice
$m$ di un operatore nello spazio degli operatori.

\rowcolors{2}{Paper}{white}
\begin{tabularx}{\textwidth}{p{30mm}>{\ttfamily\small}p{59mm}Y}
  \rowcolor{Navy}
  \color{white}\bfseries\sffamily Passo &
  \color{white}\bfseries\sffamily Nel modulo &
  \color{white}\bfseries\sffamily Significato \\
  Monopolo & T[(0,0)] = I / sqrt(d) & Traccia normalizzata \\
  Peso massimo & highest = J\_plus ** ell & Avvio della catena di rango $\ell$ \\
  Normalizzazione & highest / norm(highest) & Norma Hilbert-Schmidt unitaria \\
  Discesa & J\_minus @ T - T @ J\_minus & Generazione di $m-1$ tramite commutatore \\
\end{tabularx}

\begin{notebox}
  \cardtitle{Perch\'e questa via \`e robusta}
  Evita dipendenze dall'ordine degli argomenti, dalle normalizzazioni e dalle
  fasi adottate da specifiche implementazioni dei coefficienti di
  Clebsch-Gordan. La costruzione resta equivalente, ma \`e controllata
  direttamente dall'algebra del momento angolare.
\end{notebox}

\newpage

% -----------------------------------------------------------------------------
% 3. Diagnostica della diffusione
% -----------------------------------------------------------------------------
\section*{3. La diagnostica della diffusione}

La diffusione isotropa del modello \`e generata dalla doppia azione dei
commutatori con le tre componenti del momento angolare:

\begin{equationbox}
\[
  \frac{d\rho}{dt}
  =-D\sum_{\alpha=x,y,z}
  [J_{\alpha},[J_{\alpha},\rho]].
\]
\end{equationbox}

I tensori $T_{\ell m}$ diagonalizzano questo generatore. Il relativo autovalore
dipende soltanto dal rango $\ell$, non da $m$:

\begin{equationbox}
\[
  \sum_{\alpha=x,y,z}
  [J_{\alpha},[J_{\alpha},T_{\ell m}]]
  =\ell(\ell+1)T_{\ell m}.
\]
\end{equationbox}

Da qui segue la previsione verificata numericamente:

\begin{equationbox}
\[
  \boxed{
    c_{\ell m}(t)
    =e^{-D\ell(\ell+1)t}c_{\ell m}(0)
  }.
\]
\end{equationbox}

\subsection{Velocit\`a attese per il caso \texorpdfstring{$j=2$}{j=2}}

\rowcolors{2}{Paper}{white}
\begin{tabularx}{\textwidth}{C{17mm}C{49mm}Y}
  \rowcolor{Navy}
  \color{white}\bfseries $\ell$ &
  \color{white}\bfseries Fattore di decadimento &
  \color{white}\bfseries Lettura fisica \\
  $0$ & $1$ & La traccia resta costante \\
  $1$ & $e^{-2Dt}$ & Il dipolo sopravvive pi\`u a lungo \\
  $2$ & $e^{-6Dt}$ & Il quadrupolo si attenua pi\`u rapidamente \\
  $3$ & $e^{-12Dt}$ & Le strutture fini vengono eliminate presto \\
  $4$ & $e^{-20Dt}$ & Il massimo dettaglio accessibile \`e il pi\`u fragile \\
\end{tabularx}

\subsection{Che cosa viene controllato}

\begin{tcbraster}[
  raster columns=3,
  raster equal height=rows,
  raster column skip=3mm
]
  \begin{notebox}
    \cardtitle{Coefficienti}
    Per ogni stato della traiettoria si calcola
    $c_{\ell m}(t)=\operatorname{Tr}[\rho(t)T_{\ell m}^{\dagger}]$.
  \end{notebox}
  \begin{notebox}[borderline west={2.4pt}{0pt}{Gold}]
    \cardtitle{Potenza}
    Si aggregano i componenti nello stesso rango:
    $P_{\ell}(t)=\sum_m|c_{\ell m}(t)|^2$.
  \end{notebox}
  \begin{notebox}[borderline west={2.4pt}{0pt}{Navy}]
    \cardtitle{Errore}
    Si misura lo scarto massimo tra evoluzione osservata e previsione teorica.
  \end{notebox}
\end{tcbraster}

\vspace{4mm}
\begin{notebox}
  \cardtitle{Conclusione operativa}
  Se tutti i componenti con lo stesso $\ell$ seguono la medesima legge
  esponenziale, la dinamica numerica rispetta la simmetria rotazionale e il
  generatore teorico. Se un rango devia, la discrepanza segnala un problema
  nella costruzione degli operatori, nella vettorizzazione della supermatrice
  o nell'applicazione del canale.
\end{notebox}

\vfill
\muted{Riferimento nel progetto: \texttt{spin\_multipoles.py}. La nota descrive
la diagnostica e non propone modifiche alle ipotesi fisiche del modello.}

\end{document}
